% /root/monodromy, 2026-11-19. New complete argument; independent audits pending.
\subsection{Two boundary degree comparisons for the even images}
\label{subsec:ss-all-even-groups-candidate}
\status{Candidate pending complete independent audits. This argument uses
the separately audited first-image sheaf. It proves the missing geometric
hypothesis by two different odd-component boundary decompositions, not by
extrapolation from the retained degree-two pencil experiments.}

Use $X=E^2$, $E:y^2+y=x^3$, $L=\mathcal O(3O)\boxtimes\mathcal O(3O)$,
and the polynomial $h_d^{\mathrm u}$ of
Lemma~\ref{lem:ss-even-image-sheaf-candidate}.
Work over an algebraic closure until the final group statement.

\begin{lemma}[The intersection-point Hasse divisor]
\label{lem:ss-intersection-Hasse-degree-candidate}
Let $a,e$ be positive odd integers and let $D\in|L^e|$ be a fixed
smooth sharp curve. On $P_a=|L^a|$, the divisor $R_D$ parametrizing
nonreduced intersections $C\cap D$ is irreducible and reduced,
and has degree
\[
                       \deg R_D=18ae+9e^2.
\]
Its generic member has exactly one intersection of multiplicity
two and all other intersections simple. Its Frobenius equation
is the determinant on the finite intersection algebra.
\end{lemma}

\begin{proof}
Remove the locus of sections vanishing identically on $D$.
It is empty if $a<e$. If $a=e$ it is a point in a projective
space of dimension $9a^2-1\ge8$. If $a>e$ its codimension is
\[
                  9a^2-9(a-e)^2=9e(2a-e)>1.
\]
On the remaining open $B$, the intersection family
$Z\subset D\times B$ is finite flat of degree $n=18ae$.
Indeed it is a relative effective Cartier divisor, defined by a
section nonzero on every curve fibre; it is flat and has
zero-dimensional fibres, so properness makes it finite.
Write $\mathcal A=\pi_*\mathcal O_Z$. Cohomology of the divisor
sequence gives an exact sequence of vector bundles
\[
 0\to\mathcal O_B\to\mathcal A\to
 H^1(D,L^{-a}|_D)\otimes\mathcal O_B(-1)
       \to H^1(D,\mathcal O_D)\otimes\mathcal O_B\to0.
\]
Here $g(D)=9e^2+1$ and Riemann--Roch gives
$h^1(D,L^{-a}|_D)=n+g(D)-1=n+9e^2$.
Thus $\det\mathcal A=\mathcal O_B(-n-9e^2)$.
Its Frobenius determinant is a section of
$\mathcal O_B(n+9e^2)$ and extends over the removed
codimension-at-least-two locus. It vanishes precisely when the
finite geometric intersection scheme is nonreduced.

The double-zero incidence over $D$ is a projective bundle:
the two conditions on the value and tangent derivative at any
point are independent because $L^a$ is very ample.
It is irreducible of dimension $\dim P_a-1$.
We check the generic contact and its multiplicity.
For $a\ge2$, the two factors of degree $3a\ge6$ interpolate
first jets at two points and second jets at one point. Tensoring
these interpolations gives full joint surface first jets and
full surface second jets. This follows directly from
$H^1(E,\mathcal O_E(3aO)(-T))=0$ for every effective
subscheme of length at most four. Hence a prescribed double
zero can be made exactly double, and another double zero
imposes two further conditions with its point varying in a
curve. The latter incidence is proper inside the double-zero
incidence. If $a=1$, choose the point of $D$ away from the two
elliptic flex sets. The joint curve jets at this point and a
second point have rank four except at finitely many points and
at least three there: the exceptional surface-jet locus is the
union of two coordinate fibres and finitely many crossed
tangencies, and $D$ is integral and ample. The same dimension
count applies, and the product-cubic second jets give exact
contact two. This is the complete double-zero calculation also
used in Lemma~\ref{lem:ss-even-one-loss-point-candidate}.

Thus the generic member of the incidence has exactly one
double zero. The local length-two factor of its finite algebra
has a Weierstrass equation
\[
                             z^2+b'z+c.
\]
The value and first-jet interpolation makes $b',c$ independent
local deformation parameters. Frobenius in the basis $1,z$
has matrix $\left(\begin{smallmatrix}1&c\\0&b'\end{smallmatrix}\right)$,
so its determinant is $b'$, with vanishing order one.
All other points contribute unit determinants.
The marked double-point incidence may map generically inseparably
of degree two to its image in characteristic two; no separability
of that incidence map is used. The displayed determinant directly
computes the reduced equation on the coefficient base itself.
Consequently the extended Frobenius determinant has the reduced
irreducible double-zero divisor as its entire divisor. Its
degree is exactly $n+9e^2$, as asserted.
\end{proof}

\begin{lemma}[An exact factorization on an odd-component boundary]
\label{lem:ss-even-boundary-factorization-candidate}
Let $a\ge3$ and $e\ge1$ be odd, put $d=a+e$, and fix a smooth
sharp $D\in|L^e|$. The linear embedding
\[
 j_D:P_a\hookrightarrow P_d,\qquad C\longmapsto C\cup D
\]
satisfies, up to a nonzero scalar,
\[
                         j_D^*h_d^{\mathrm u}=h_a\,r_D,
\]
where $h_a$ is the irreducible reduced odd-degree relative
Hasse polynomial and $r_D$ is the irreducible equation of
$R_D$ from the preceding lemma.
For $a=1$, the same restriction is nonzero and is divisible
by $r_D$; no complete factorization in this case is needed below.
\end{lemma}

\begin{proof}
A generic smooth sharp $C$ meets $D$ transversely. Its nodal
union has normalization defect $2+2=4$, so its stable rank is
$9d^2-3$, the generic even rank. At each transverse node the
mixed derivative $\delta_1\delta_2(s_Cs_D)$ is nonzero, as in
Lemma~\ref{lem:ss-even-one-loss-point-candidate}. Thus the union
lies in the vector-bundle domain $P^\circ$ and the restricted
polynomial is not zero.

For $a\ge3$, the odd Hasse divisor is irreducible and reduced
of degree $9a^2$, by Theorem~\ref{thm:ss-all-odd-Hasse}.
Its generic curve is smooth. To check this explicitly, the
singular-curve discriminant for $|L^a|$ is irreducible of
degree $c_2(J^1L^a)=54a^2$: joint first jets and second jets
give a generic unique ordinary node, with separable degree one
for the singular-point incidence. Its degree differs from
$9a^2$, so it cannot equal the Hasse divisor.
Also $R_D$ has degree $18ae+9e^2\ne9a^2$; equality would give
$(a+e)^2=2a^2$, impossible for positive integers.
Thus the generic curve on $h_a=0$ is smooth and transverse to
$D$. Its defect is at least three, so the union defect is at
least five and its unit-image determinant vanishes.
There remain transverse intersection nodes, so these points
are in $P^\circ$. Hence $h_a$ divides $j_D^*h_d^{\mathrm u}$.

The generic member of $R_D$ is smooth and sharp as a source
curve. For $a\ge3$ this follows from irreducibility and the
degree comparisons with $h_a$ and the singular-curve
discriminant. The latter degree cannot equal $18ae+9e^2$,
since equality would give $(a+e)^2=7a^2$.
For $a=1$, the tangent-six-space construction in
Lemma~\ref{lem:ss-even-one-loss-point-candidate} supplies a
smooth sharp member directly, so the same assertion follows
by irreducibility of $R_D$.
Its union with $D$ has one simple cross tacnode and otherwise
nodes. Its defect is $2+2+1=5$, so the restricted unit-image
Hasse polynomial vanishes on $R_D$. The remaining ordinary
nodes again place the generic point in $P^\circ$.
This proves the claimed divisibility by $r_D$ in all cases.

For $a\ge3$ the two irreducible factors are distinct, and
their degrees add to the full degree of the restriction:
\[
             9a^2+(18ae+9e^2)=9(a+e)^2=9d^2.
\]
The nonzero restriction therefore equals their product, with
both multiplicities exactly one. No unproved comparison of
coherent first-image lattices is used: divisibility follows
from stable ranks, and the equality follows from this degree
identity.
\end{proof}

\begin{theorem}[The even unit-image Hasse divisor]
\label{thm:ss-even-unit-Hasse-irreducible-candidate}
For every even $d\ge4$, the divisor $h_d^{\mathrm u}=0$ is
geometrically irreducible and reduced.
\end{theorem}

\begin{proof}
Use the boundary with $(a,e)=(d-1,1)$. Its restricted
polynomial is a product of two distinct irreducible polynomials
of degrees
\[
                        9(d-1)^2,\qquad18d-9.
\]
Every nonzero homogeneous global factor restricts to a nonzero
homogeneous polynomial of the same positive degree: the whole
restriction is nonzero, and $j_D$ is a linear projective
embedding. The squarefree restriction thus proves that
$h_d^{\mathrm u}$ is reduced and has at most two irreducible
components. If there are two, their degrees must be exactly
the displayed pair.

For $d\ge6$, use instead a fixed sharp degree-three component,
so $(a,e)=(d-3,3)$. The same argument would force the degrees
of the same two global components to be
\[
                       9(d-3)^2,\qquad54d-81.
\]
The two unordered pairs are different. Equality in the same
order would give $(d-1)^2=(d-3)^2$, impossible for $d\ge6$.
Equality in the opposite order would give
\[
                  d^2-8d+10=0,
                  \quad\text{or}\quad(d-4)^2=6,
\]
also impossible for an integer. This excludes two components.

For $d=4$, the first pair is $81,63$.
Fix a sharp degree-three curve and vary a degree-one curve.
The preceding lemma shows that the restriction of $h_4^{\mathrm u}$
is nonzero and divisible by the irreducible point-intersection
factor of degree $18\cdot1\cdot3+9\cdot3^2=135$.
If the global divisor had two components of degrees $81,63$,
their nonzero restrictions would have those same degrees.
Their product cannot be divisible by an irreducible polynomial
of degree $135$, since such a polynomial must divide one of
the two factors. This is the remaining contradiction.
\end{proof}

\begin{theorem}[All degrees on the fixed supersingular product]
\label{thm:ss-product-all-full-groups-candidate}
Let $k$ be any perfect field of characteristic two, and embed
$X=E\times E$, $E:y^2+y=x^3$, by
$L=\mathcal O_E(3O)\boxtimes\mathcal O_E(3O)$.
For every $d\ge1$, the entire geometric and arithmetic images
of the relative mod-two Hasse--Witt local system are
\[
 G_{\mathrm{geom},d}=G_{\mathrm{arith},d}=
 \begin{cases}
 \GL_{9d^2-1}(\mathbf F_2),&d\text{ odd},\\[2pt]
 \mathbf F_2^{\,9d^2-4}\rtimes\GL_{9d^2-4}(\mathbf F_2),
                                      &d\text{ even}.
 \end{cases}
\]
The even image is the full pointwise stabilizer of the unique
nonzero invariant covector in rank $9d^2-3$. The odd system
is geometrically irreducible; the even system has its unique
proper nonzero invariant subspace of codimension one.
\end{theorem}

\begin{proof}
The odd degrees are Theorem~\ref{thm:ss-all-odd-full-groups}.
For even $d\ge4$, apply the preceding Hasse theorem to the
conditional reduction in
Theorem~\ref{thm:ss-even-full-group-if-unit-Hasse-irreducible-candidate}.
Degree two is the separately proved exact certificate
Theorem~\ref{thm:ss-product-degree-two-full-group}.
The natural translation module and general linear block of its
full stabilizer show that its common covector kernel is also
the unique proper nonzero invariant subspace.
All geometric and arithmetic identifications, and the passage
to ambient hypersurfaces, are those already checked in these
theorems. No claim for an arbitrary ambient variety follows
from this fixed-surface classification.
\end{proof}
